
\documentclass[12pt]{article}
\usepackage{amsmath, amssymb}
\usepackage{geometry}
\geometry{margin=1in}
\title{SAT.O2 – Emergent Gravitational Action}
\author{SATO.4D Coherence Project}
\date{June 2025}

\begin{document}
\maketitle

\section*{1. Foundational Assumptions}

\begin{itemize}
    \item Spacetime is a 4D manifold \( M \) populated by 1D physical filaments \( \gamma: \mathbb{R} \to M \).
    \item No metric, connection, or field is assumed a priori.
    \item All geometric structure — including curvature — emerges statistically from the filament ensemble.
    \item The resolving surface \( \Sigma_t \) provides the time foliation field \( u^\mu(x) \).
\end{itemize}

\section*{2. Filament Ensemble and Local Geometry}

Let \( F(x) \) denote the set of filaments intersecting a point \( x \in M \). Define the tangent vector of filament \( \gamma \) at affine parameter \( \lambda \):
\[
v^\mu(\lambda) = \frac{d\gamma^\mu}{d\lambda}
\]

\section*{3. Emergent Metric}

Define the emergent co-metric via ensemble averaging:
\[
\tilde{g}_{\mu\nu}(x) = \left\langle v_\mu v_\nu \right\rangle_{F(x)}
\]
Assuming statistical isotropy and sufficient density, the inverse metric \( g^{\mu\nu}(x) \) exists:
\[
g_{\mu\nu}(x) = \left( \tilde{g}_{\mu\nu}(x) \right)^{-1}
\]

\section*{4. Emergent Connection}

The Levi-Civita connection emerges from the metric via:
\[
\Gamma^\lambda_{\mu\nu}(x) = \frac{1}{2} g^{\lambda\rho} \left( \partial_\mu g_{\rho\nu} + \partial_\nu g_{\rho\mu} - \partial_\rho g_{\mu\nu} \right)
\]
This connection is:
\begin{itemize}
    \item Torsion-free,
    \item Metric-compatible: \( \nabla_\lambda g_{\mu\nu} = 0 \)
\end{itemize}

\section*{5. Curvature Tensor}

Curvature emerges from distortion of the filament congruence:
\[
R^\rho_{\ \sigma\mu\nu}(x) = \partial_\mu \Gamma^\rho_{\nu\sigma} - \partial_\nu \Gamma^\rho_{\mu\sigma} + \Gamma^\lambda_{\nu\sigma} \Gamma^\rho_{\mu\lambda} - \Gamma^\lambda_{\mu\sigma} \Gamma^\rho_{\nu\lambda}
\]

\subsection*{5.1 Ricci Tensor and Scalar}

\[
R_{\mu\nu}(x) = R^\rho_{\ \mu\rho\nu}, \quad R(x) = g^{\mu\nu} R_{\mu\nu}
\]

\section*{6. Emergent Stress-Energy Tensor}

From local filament energy:

\begin{itemize}
    \item \( E_{\text{vib}}(x) \): Vibrational energy density,
    \item \( L_{\text{link}}(x) \): Topological linking density
\end{itemize}

Define:
\[
T_{\mu\nu}(x) = \left\langle E_{\text{vib}}(x), L_{\text{link}}(x) \right\rangle_{F(x)}
\]

\section*{7. Emergent Gravitational Field Equations}

\[
\left\langle R_{\mu\nu} - \frac{1}{2} g_{\mu\nu} R \right\rangle_{F} = \kappa \, T_{\mu\nu}(x)
\]
with:
\[
\kappa = \frac{8\pi G}{c^4}
\]
These are the emergent Einstein Field Equations in SAT4D: curvature arises statistically from filament strain and alignment.

\section*{8. Structural Interpretation}

\begin{itemize}
    \item Curvature is a measure of ensemble distortion, not a pre-existing field.
    \item The strain tensor \( S_{\mu\nu} = \nabla_\mu u_\nu + \nabla_\nu u_\mu \) links directly to emergent curvature.
    \item \(\theta_4(x)\) may be used to assess local inertial deviation, but is not a driver of curvature.
\end{itemize}

\section*{9. Deviation from Classical GR}

SAT predicts:
\begin{itemize}
    \item Deviations from Einstein gravity in regions of sparse or asymmetric filament density.
    \item Curvature singularities avoided due to topological regularization.
    \item Falsifiability via correlation between \( T_{\mu\nu}(x) \) and measurable linking density.
\end{itemize}

\section*{10. Frame Declaration}

This derivation is conducted in \textbf{Interpretive Mode 1 (True Block)}. All dynamics emerge from slicing structure; no internal filament evolution is assumed.

\end{document}
